%% =========================================================================
%% chapters/ch13-social-mechanics.tex
%%
%% Chapter 13. Infection between hosts, built on the definition of disease
%% derived in Chapter 1: a pathogen is a lesion of kind (L2) -- an addition
%% of reactions to a host's network -- and transmission is the process by
%% which that lesion is delivered. The chapter derives what contact
%% structure contributes to that delivery, and separates it carefully from
%% what the environment contributes, which it does not attempt to deny.
%%
%% This chapter replaces an earlier one built on template matching, whose
%% central claims it corrects. See section "What this replaces".
%%
%% Every link carries a label: [D] derived here, [S] supplied with a
%% reference, [C] checked here against a computation (validation/ch13/).
%%
%% Labels in this chapter are prefixed sm: .
%% =========================================================================

\chapter{Transmission Between Hosts}
\label{ch:social-mechanics}
\label{sm:ch:social-mechanics}

\begin{quote}
\itshape
It is not the seed, it is the soil.
\upshape\medskip\\
\hfill --- attributed, in various forms, to nineteenth-century hygienists
\end{quote}

%% -------------------------------------------------------------------------
\section{What this chapter is about}
\label{sm:sec:intro}
%% -------------------------------------------------------------------------

\Cref{ch:what-disease-is} defined disease inside one system. A lesion is a
change to a cell's reaction network of one of three kinds: a reaction
deleted, a reaction added, or a conductance changed. A lesion is
pathogenic when it narrows the set of media in which the system can keep
itself alive (\Cref{c1:def-disease}). That definition is complete for an
inherited enzymopathy, where the lesion is already present and the question
is only what it costs.

An infection is different in one respect: the lesion arrives. A pathogen is
a lesion of kind (L2), an addition of reactions to the host's network, and
it is delivered from another host. So an account of infectious disease needs
two things that \Cref{ch:what-disease-is} does not supply. It needs to say
what kind of lesion a pathogen is, which is a statement about one host, and
it needs to say what governs delivery, which is a statement about a
population. This chapter does both, in that order.

The population half is where this chapter has something to correct. In
earlier work I argued that the rate at which a pathogen reaches hosts is
governed by the \emph{coordination regime} of the host population, and that
environmental variables---temperature, humidity, rainfall---do not enter the
calculation at all. The first claim is right, and this chapter makes it
exact and checks it. The second claim is wrong, and
\Cref{sm:sec:environment} says so plainly and explains how I came to hold
it. The chapter that results is narrower than the one it replaces and, I
think, worth more for being narrower.

\subsection*{How the chapter argues}

As in \Cref{ch:what-disease-is}, each statement carries a label:
\begin{itemize}[nosep]
\item \lnkD{} \emph{derived} here from the definitions of
      \Cref{ch:what-disease-is} and this chapter;
\item \lnkS{} \emph{supplied} from the published literature, with a
      reference;
\item \lnkC{} \emph{checked} here by computation, with the code in
      \texttt{validation/ch13/}.
\end{itemize}
The \lnkC{} labels in this chapter are weaker than those of
\Cref{ch:what-disease-is} and it is worth saying why at the outset. There,
a check ran a stoichiometric computation over a network assembled from
public databases, so the thing checked was a property of the real red cell.
Here, the checks are computations on contact structure with a stipulated
per-contact transmission probability. They verify that a mathematical claim
about contact structure holds, and they verify nothing about any particular
disease anywhere. No number in this chapter is calibrated against observed
incidence. \Cref{sm:sec:ledger} collects the labels.

%% -------------------------------------------------------------------------
\section{A pathogen is an added-reaction lesion}
\label{sm:sec:pathogen-lesion}
%% -------------------------------------------------------------------------

Begin inside a single host, where \Cref{ch:what-disease-is} already applies.

\begin{definition}[Infection]
\label{sm:def-infection}
Let $\mathcal N=(N,\kappa)$ be a host system with exchanged species
$\mathcal X$ and admissible media $\mathcal M$
(\Cref{c1:def-kinetic}). An \emph{infectious agent} is a set of reactions
$P$, together with any species they introduce, that is added to
$\mathcal N$ to give $\mathcal N'=(N\cup P,\kappa')$. The addition is a
lesion of kind (L2) in the sense of \Cref{c1:def-lesion}. The agent is
\emph{pathogenic in} $\mathcal M$ if that lesion is pathogenic:
$\mathcal V(\mathcal N')\subsetneq\mathcal V(\mathcal N)$.
\end{definition}

Three consequences follow from \Cref{c1:thm-internal,c1:thm-medium,c1:thm-lesion}
of \Cref{ch:what-disease-is}, and each replaces something the earlier
chapter asserted.

\begin{proposition}[What an infection cannot change]
\label{sm:prop-no-holonomy}
\lnkD{} An infection does not change the holonomy of any cycle that the
host network retains. Internal cycles have holonomy zero before and after
(\Cref{c1:thm-internal}); a cycle through the medium has holonomy
$\langle\mu_{\mathcal X},N_{\mathcal X}z\rangle$ before and after
(\Cref{c1:thm-medium}). Any new emergent cycle the pathogen creates has its
holonomy fixed by the medium, not by the pathogen.
\end{proposition}

\begin{proof}
The proposition is \Cref{c1:thm-lesion}(a) and (d) applied to the addition
$P$, together with the observation that adding reactions does not remove
cycles.
\end{proof}

So a pathogen does not bring energy with it, and it does not alter the
thermodynamic bookkeeping of the host. What it can do is add cycles that
draw on the holonomy the host's medium already supplies. That is the precise
sense in which a pathogen is a parasite: by \Cref{c1:eq-power} the power
dissipated at steady state is drawn entirely from the medium, and a new
emergent cycle takes a share of it.

\begin{definition}[Virulence]
\label{sm:def-virulence}
The \emph{virulence} of an infectious agent $P$ in a host $\mathcal N$ and
medium $\mu_{\mathcal X}$ is the severity of the lesion it constitutes, the
loss of margin $h(\mathcal N,\mu_{\mathcal X})-h(\mathcal N',
\mu_{\mathcal X})$ of \Cref{c1:def-disease}.
\end{definition}

\begin{proposition}[Virulence is a relation, not a property]
\label{sm:prop-virulence}
\lnkD{} Virulence in the sense of \Cref{sm:def-virulence} is a function of
three arguments---the agent, the host network and the medium---and is not
determined by the agent alone. The same $P$ can have zero severity in one
host-medium pair and positive severity in another.
\end{proposition}

\begin{proof}
Severity is defined from $h(\mathcal N,\cdot)$ and $h(\mathcal N',\cdot)$,
both of which depend on $\mathcal N$ and on the medium through
$\mathcal V$ (\Cref{c1:def-viability}). For an explicit pair, let $P$ be a
single reaction consuming an exchanged species the host does not otherwise
need. In a medium holding that species at high potential the host loses
nothing. In a medium where it is scarce and the host does need it, the host
loses margin.
\end{proof}

That virulence is not an intrinsic property of a pathogen is the one claim
of the earlier chapter that survives untouched and is strengthened. There it
followed from a posited overlap integral. Here it follows from the
definition of severity, which is computed from a network and a medium.

\begin{proposition}[Why a host range is narrow]
\label{sm:prop-host-range}
\lnkD{} Let every reaction of $P$ require a species or catalytic context
absent from a host $\mathcal N$. Then $P$ adds no reaction that can carry
flux in $\mathcal N$: by \Cref{c1:prop-flux} such reactions lie on no cycle
of $\ker N_{\mathcal I}$, so every steady state of $\mathcal N$ remains a
steady state of $\mathcal N\cup P$. By \Cref{c1:thm-lesion}(c) the
viability set is then unchanged when steady states are unique, and the agent
is not pathogenic in any medium of $\mathcal M$.
\end{proposition}

\begin{proof}
A reaction whose substrate set is not a subset of the host's species has
every rate identically zero, so it contributes a column that no steady flux
can use, and the rank test of \Cref{c1:prop-flux} excludes it from every
cycle. The rest is \Cref{c1:thm-lesion}(c).
\end{proof}

\Cref{sm:prop-host-range} is the corrected form of the earlier chapter's
argument about bacteriophages. That argument was right in its conclusion and
confused in its reasoning: it said the overlap integral between the phage
template and the human host state is zero, which treats the barrier as a
failure of geometric match. The real statement needs no template. A phage
encodes reactions requiring a bacterial ribosome, a bacterial polymerase and
a bacterial envelope; a human cell has none of these; so the added
reactions cannot carry flux, and the addition is not a lesion at all. The
absence of those reactions from a human cell is checkable in principle from
annotation, in the way \Cref{c1:sec-erythrocyte} checked which erythrocyte
enzymes lie on a cycle, and I have not done it here.

\begin{remark}[What this does not explain]
\label{sm:rem-spillover}
\Cref{sm:prop-host-range} is a statement about the extreme case, where
\emph{no} reaction of the agent can run. It says nothing about the
interesting case, a zoonosis that runs poorly in a new host and then adapts.
That case needs the kinetics $\kappa$ which \Cref{c1:in-steady} assumes and
does not compute, together with a model of variation under selection, and
neither this chapter nor \Cref{ch:what-disease-is} supplies them.
\end{remark}

%% -------------------------------------------------------------------------
\section{Delivery: the population half}
\label{sm:sec:delivery}
%% -------------------------------------------------------------------------

\Cref{sm:def-infection} describes a lesion that has already arrived.
Whether it arrives is a question about a population of hosts and the
contacts between them, and this is where the coordination argument of the
earlier chapter belongs. Reduced to its defensible core, the argument is
about a second moment.

\begin{definition}[Contact structure]
\label{sm:def-contact}
A \emph{contact structure} on a population of $n$ hosts assigns to each host
$i$ a number of contacts $k_i$ per unit time, together with a rule for whom
those contacts are with. The \emph{contact distribution} is the empirical
distribution of $(k_i)$, with mean $\langle k\rangle$ and second moment
$\langle k^2\rangle$. A \emph{per-contact transmissibility} $\tau\in[0,1]$
is the probability that one contact between an infectious and a susceptible
host delivers the lesion.
\end{definition}

The quantity that decides whether an introduced lesion spreads is the
reproduction number: the expected number of hosts that one infectious host
infects in an otherwise susceptible population. For a population mixing at
random subject to its contact distribution, it is not $\tau\langle k\rangle$.

\begin{theorem}[The threshold depends on the second moment]
\label{sm:thm-second-moment}
\lnkS{} For a configuration-model contact network with contact distribution
$(k_i)$ and per-contact transmissibility $\tau$, the reproduction number is
\begin{equation}
R_0=\tau\,\frac{\langle k^2\rangle}{\langle k\rangle},
\label{sm:eq-r0}
\end{equation}
and the epidemic threshold is at $R_0=1$. Equivalently, for a population
divided into groups of activity $k_i$ mixing proportionately, with contact
matrix $C_{ij}=k_ik_j/(g\langle k\rangle)$, the dominant eigenvalue of the
next-generation matrix $\tau C$ is $\tau\langle k^2\rangle/\langle
k\rangle$.
\end{theorem}

\Cref{sm:eq-r0} is standard~\citep{anderson1991infectious,diekmann1990host},
and I supply rather than derive it. The mechanism is that a contact is more
likely to lead to a host of high degree than a uniformly chosen host, by a
factor proportional to that degree. Two consequences do the work of this chapter.

\begin{corollary}[Mean contact does not determine vulnerability]
\label{sm:cor-mean}
\lnkD{} Two populations with the same mean contact number
$\langle k\rangle$ and the same $\tau$ can have reproduction numbers
differing by the ratio of their $\langle k^2\rangle/\langle k\rangle$, which
is unbounded. In particular a population can be below threshold while
another with identical mean contact and an identical pathogen is far above
it.
\end{corollary}

\begin{proof}
Immediate from \Cref{sm:eq-r0}, since $\langle k^2\rangle$ is not determined
by $\langle k\rangle$. For unboundedness, put a fraction $1/m$ of the
population at degree $m\langle k\rangle$ and the rest at degree near zero.
\end{proof}

\begin{corollary}[Structure can move the threshold with no pathogen change]
\label{sm:cor-structure}
\lnkD{} In \Cref{sm:eq-r0} the pathogen enters only through $\tau$ and the
environment does not enter at all. Hence any change in the contact
distribution that raises $\langle k^2\rangle/\langle k\rangle$ raises $R_0$
at fixed $\tau$, and a sufficiently large such change crosses the threshold
without any change in the pathogen, the hosts' networks, or the
surroundings.
\end{corollary}

\paragraph{Check: same mean, different threshold.} \lnkC{} Hold mean contact at
exactly $8$ and $\tau$ at $0.08$. Three contact distributions then give
$\langle k^2\rangle/\langle k\rangle$ of $8.0$ (every host identical),
$9.0$ (Poisson) and $49.9$ (a tenth of hosts with ten times the contacts),
hence $R_0$ of $0.64$, $0.72$ and $3.99$. Bond percolation on sampled
graphs with those degree sequences gives mean outbreak fractions of
$0.003$, $0.004$ and $0.203$: the first two die out, the third does not.
The ratio of reproduction numbers is about $6$ at identical mean contact.
Code: \texttt{validation/ch13/coordination.py}, check C1.

\paragraph{Check: the matrix form.} \lnkC{} Letting five groups differ in
activity at constant population mean contact, under proportionate mixing,
raises the dominant eigenvalue of the next-generation matrix from $0.500$
to $0.703$, matching $\tau\langle k^2\rangle/\langle k\rangle$ to the
printed digits. Check C2b.

\subsection{Group size, and why concentration is self-limiting}

Contact is not only unequal, it is grouped: households, wards, barracks,
markets. Grouping matters through a mechanism \Cref{sm:eq-r0} does not
capture, because within a group the same hosts are contacted repeatedly and
transmission compounds over generations.

\begin{definition}[Household reproduction number]
\label{sm:def-rstar}
Partition the population into groups of size $s$ and let each host spend
$s-1$ of its contacts inside its group and the remaining
$\langle k\rangle-(s-1)$ outside. Let $\mu(s,\tau)$ be the expected number
of hosts ultimately infected within a group following one introduction into
it. The \emph{household reproduction number} is
\begin{equation}
R_*=\mu(s,\tau)\,\tau\,\bigl(\langle k\rangle-(s-1)\bigr).
\label{sm:eq-rstar}
\end{equation}
\end{definition}

\begin{proposition}[Concentration is non-monotone]
\label{sm:prop-nonmonotone}
\lnkD{} $R_*$ of \Cref{sm:eq-rstar} is not monotone in $s$ at fixed
$\langle k\rangle$. The factor $\mu$ is non-decreasing in $s$ and rises
sharply once $(s-1)\tau$ approaches $1$, while the factor
$\langle k\rangle-(s-1)$ falls linearly and reaches zero at
$s-1=\langle k\rangle$, where $R_*=0$.
\end{proposition}

\begin{proof}
$\mu$ is non-decreasing because a larger group has more susceptibles
reachable by the same chain. The second factor is affine and decreasing in
$s$, vanishing at $s=\langle k\rangle+1$, at which point the product
vanishes while $\mu>1$.
\end{proof}

\paragraph{Check: group size at fixed mean contact.} \lnkC{} Hold mean contact at
$40$ and $\tau$ at $0.05$, and treat within-group spread as a Reed--Frost
chain compounded to extinction. Then $R_*$ rises from $2.00$ at $s=1$ to
$6.19$ at $s=30$, while the within-group final size $\mu$ rises from $1.00$
to $11.25$ and the outside budget falls from $40$ to $11$. Over this range
the amplification dominates. Check C3.

\Cref{sm:prop-nonmonotone} is worth dwelling on because it is the kind of
conclusion the earlier chapter could not reach. There, vulnerability
increased monotonically with coordination by construction, since the
coordination weighting $f$ was a non-decreasing step function. Here the
competing effects are visible: concentrating contact amplifies locally and
starves transmission between groups, and which wins depends on the numbers.
A population that has concentrated \emph{all} of its contact into closed
groups transmits nothing between them. That is why quarantine works, and the
monotone account cannot say so.

%% -------------------------------------------------------------------------
\section{Two claims of the earlier chapter that fail}
\label{sm:sec:failures}
%% -------------------------------------------------------------------------

\subsection{The coordination weighting function}
\label{sm:sec:weighting}

The earlier chapter parametrised a population by a scalar order parameter
$r\in[0,1]$, cut $[0,1]$ into five named regimes, and assigned each a
vulnerability weight: $f=0.2$ for turbulent ($r<0.3$) rising in steps to
$f=1.0$ for phase-locked ($r\ge0.95$). The ratio $f(0.95)/f(0.3)=5$ was
reported as a theorem: a phase-locked population is five times as vulnerable
as a turbulent one.

By \Cref{sm:eq-r0} the structural quantity this $f$ is standing in for is
$\langle k^2\rangle/\langle k\rangle$. So the function is testable: pick a
family of contact distributions indexed by $r$ at fixed mean, and compute.

\paragraph{Check: the weighting function is not monotone.} \lnkC{} Take a
two-level family in which a share $r$ of all contacts is absorbed by a core
holding a share $1-r$ of the population, at mean contact fixed at $8$. Then
$\langle k^2\rangle/\langle k\rangle$ takes the values $38.7$, $9.3$,
$11.2$, $46.7$, $250.9$ at $r=0.15,0.40,0.65,0.87,0.97$. The function is
\emph{U-shaped}, with its minimum near $r=0.4$, not increasing. Check C4.

The reason is visible once stated. At fixed mean, $\langle k^2\rangle$ is
minimised by the uniform distribution and rises in \emph{both} directions
away from it. At $r=0.15$ a small periphery carries degree $45$ while $85\%$
of the population carries $1.4$; at $r=0.97$ a $3\%$ core carries degree
$259$. Both are concentrated; concentration is what the second moment
measures, and a scalar running from $0$ to $1$ does not say which tail is
heavy. \lnkD

\begin{remark}[What the weighting function was tracking]
\label{sm:rem-weighting}
The monotone intuition behind $f$ is not empty: it tracks the real fact of
\Cref{sm:cor-structure}, that concentrating contact raises the threshold
quantity. It fails as a function of $r$ because $r$ was never defined so as
to determine the contact distribution. I therefore drop $f$ and its
numerical values entirely, and keep \Cref{sm:eq-r0}, which needs no regime
names and is exact.
\end{remark}

\subsection{Independence from the environment}
\label{sm:sec:environment}

The earlier chapter made a stronger claim in three parts. Disease prevalence
does not depend on temperature, humidity, rainfall or season, except
insofar as these change host states or contact structure. Environmental
variables ``do not appear in the prevalence equation''. A tropical region
made artificially cool would show unchanged prevalence. I withdraw all
three. They are false for an important class of diseases, and the way they
are false is instructive.

The argument that produced it was formally valid and materially
unsound. \Cref{sm:eq-r0} indeed contains no temperature. But $\tau$ is not
a constant of the pathogen; it is the probability that a contact transmits,
and for any disease with a vector or an environmental stage, that
probability is itself a function of the environment. Writing
$\tau=\tau(\text{environment})$ is not a loophole in the framework, it is
the framework applied honestly: the chapter had placed the environment
outside the equation by declaring $\tau$ fixed, and then discovered it was
not there.

For malaria the dependence is measured and sharp. Transmission of
\emph{Plasmodium falciparum} by \emph{Anopheles} is a thermal performance
curve built from the separate temperature responses of mosquito survival,
biting rate and the parasite's development time. A model assembled from
those responses puts the optimum at $25\,^{\circ}$C and predicts that
transmission declines steeply above $28\,^{\circ}$C, and field measurements
of the entomological inoculation rate across Africa support both the
optimum and the decline~\citep{mordecai2013optimal}. \lnkS{} A later
re-estimation across eleven mosquito-borne pathogens puts the optimum for
\emph{P.\ falciparum} in \emph{An.\ gambiae} at $25.4\,^{\circ}$C, with
transmission falling to zero below about $19\,^{\circ}$C and above about
$33\,^{\circ}$C~\citep{mordecai2019thermal}. \lnkS

The geographic consequence has been mapped. A degree-day model of parasite
development, requiring a minimum temperature of $16\,^{\circ}$C for
\emph{P.\ falciparum}, identifies the regions where temperature precludes
transmission throughout the year. The boundaries of suitability it finds
follow latitudinal and continental temperature gradients, punctuated by
highland areas: the Andean ridge, the tropical highlands of East Africa,
the Indonesian archipelago~\citep{gething2011modelling}. \lnkS{} So a
tropical region made artificially cool is not a thought experiment; it is a
highland.

Not every highland is excluded, and the same work finds western Kenyan
highlands where temperature precludes transmission at no point in the year.
But the earlier chapter predicted that cooling a tropical region leaves
prevalence unchanged, and that prediction is contradicted by a mapped
temperature constraint whose stated purpose is to mark where transmission
cannot occur.

\begin{proposition}[Where the environment enters]
\label{sm:prop-environment}
\lnkD{} In \Cref{sm:eq-r0} the environment can enter at exactly three
places. It can enter through $\tau$, if transmission has a stage outside a
host; through the contact distribution, if behaviour responds to conditions;
or through the host networks and media of \Cref{sm:def-infection}, which fix
whether an arriving lesion is pathogenic at all. Only the second was
represented in the earlier chapter.
\end{proposition}

\begin{proof}
$R_0$ is a function of $\tau$ and the contact distribution alone, and
whether an arriving agent is pathogenic is \Cref{sm:def-infection}; so any
environmental dependence must act through one of the three.
\end{proof}

So the corrected position is not that the environment is irrelevant, nor
that it is paramount, but that it enters at identifiable places and differs
by disease. For a directly transmitted respiratory virus, the dominant
environmental route is plausibly behavioural---through the contact
distribution---and the earlier chapter's instinct was reasonable. For a
vector-borne or water-borne disease, the environment acts on $\tau$ directly
and is not reducible to behaviour. A framework that cannot tell these apart
is not simpler; it is less informative.

\begin{remark}[Seasonality]
\label{sm:rem-seasonality}
The earlier chapter explained influenza's winter peak entirely by indoor
crowding, that is, by the contact distribution. Indoor contact is a real
contributor and the mechanism is the one in
\Cref{sm:cor-structure}. Whether it is the dominant one is an empirical
question about which the literature is not settled, and I no longer claim it
is. I state the mechanism and not the attribution.
\end{remark}

%% -------------------------------------------------------------------------
\section{Co-evolution, and what it can and cannot explain}
\label{sm:sec:coevolution}
%% -------------------------------------------------------------------------

The last structural claim of the earlier chapter was that pathogen
populations refine themselves against their hosts over time, so that regions
with longer host occupancy carry better-matched pathogens and hence more
disease. The mechanism is real and standard; the application was not
supported.

\begin{inputstmt}[Selection on an added lesion]
\label{sm:in-selection}
\lnkS{} A population of infectious agents varying heritably in the
reactions they add to a host, and differing thereby in the number of
secondary infections they cause, changes in composition toward the variants
causing more~\citep{fisher1930genetical,maynard1982evolution}. For
host-pathogen systems specifically, the direction of change in virulence
depends on how virulence trades off against transmission, and is not
toward maximum harm~\citep{anderson1979regulation,ewald1994evolution}.
\end{inputstmt}

\begin{proposition}[Co-evolution acts on the kinetics, not the holonomy]
\label{sm:prop-coev}
\lnkD{} Under \Cref{sm:in-selection}, what selection can change in the
sense of \Cref{c1:def-lesion} is which reactions the agent adds, and their
conductances. By \Cref{sm:prop-no-holonomy} it cannot change the holonomy of
any cycle. So an adapted pathogen is one whose added reactions carry more
current in the hosts it encounters, not one that has found a larger
thermodynamic driving force.
\end{proposition}

\begin{proof}
Selection changes the agent's genome and hence $P$ and $\kappa'$, which are
the arguments of \Cref{sm:def-infection}; \Cref{sm:prop-no-holonomy} fixes
the holonomies independently of both.
\end{proof}

\Cref{sm:prop-coev} is worth having, and it is as far as the argument goes.
The earlier
chapter went considerably further: it gave a settlement-history ratio of
$500{,}000$ to $50{,}000$ years for tropical against temperate human
occupation, inferred a tenfold co-evolution-time ratio, and presented
regional prevalence figures as confirmation. Those figures came from a
saturation model chosen for the purpose, not from epidemiological data; the
agreement they displayed was arithmetic. I drop the quantitative claim.
What remains is \Cref{sm:prop-coev} and the direction of the effect, with no
ratio attached.

\begin{remark}[Latency]
\label{sm:rem-latency}
The earlier chapter gave an account of tuberculosis latency: infection in
childhood, reactivation decades later when the host state crosses a
threshold. Its shape is right in the present framework. It describes an
added lesion whose severity is near zero in one host state and positive in
another, which is \Cref{c1:cor-medium} with the host's own trajectory in
place of the medium's. But the mechanism it offered for the threshold crossing, a
respiratory physiology adapted to lifelong exposure to fire smoke, is an
assumption of that framework and not an established finding.

Its scale has been estimated. Latent infection was estimated
at $23.0\%$ of the world's population in 2014, about $1.7$~billion
people~\citep{houben2016global}. \lnkS{} The lifetime risk that an
individual with latent infection progresses to active disease is estimated
at $5$ to $10$ per cent, usually within the first five years after
infection~\citep{who2018ltbi}. \lnkS{} The mechanism of late reactivation I
leave open. Ageing does bring a general increase in vulnerability to
infectious disease, through the progressive deterioration of innate and
adaptive immune responses often called
immunosenescence~\citep{goronzy2013understanding,nikolichzugich2018twilight}.
\lnkS{} But those reviews are not about tuberculosis, and I have found no
source that settles the reactivation mechanism; one line of work argues for
inflammation rather than immune decline. So I state the phenomenon, name the
domain the explanation is likely to come from, and claim neither the
fire-circle mechanism nor a specific immunological one.
\end{remark}

%% -------------------------------------------------------------------------
\section{Age structure}
\label{sm:sec:age}
%% -------------------------------------------------------------------------

Contact is strongly structured by age, and this is measured rather than
posited: diary studies across European populations find pronounced
age-assortative mixing, with the heaviest contact rates among
school-age children and between them and their
parents~\citep{mossong2008social}. \lnkS{} In the terms of this chapter, the
contact matrix $C_{ij}$ of \Cref{sm:thm-second-moment} has large diagonal
blocks at young ages.

Two cautions follow from the checks, and they cut against an easy
conclusion.

First, by check C2a, assortativity \emph{by itself} does not move the
threshold. \lnkC{} If every group's total contact is held fixed and the
groups are symmetric, then reallocating contact between within-group and
between-group leaves the dominant eigenvalue of the next-generation matrix
exactly unchanged, at $0.500$ in every case computed. The reason is
Perron--Frobenius: a non-negative matrix with constant row sums has
dominant eigenvalue equal to that row sum, with the uniform vector as
eigenvector, however the mass in each row is arranged. \lnkD{} This check
was written expecting a rise and found none.

Second, what does move the threshold is inequality in the row sums
themselves---in how much contact each group has, not with whom
(check C2b). So the epidemiologically consequential fact in the contact data
is that children have \emph{more} contacts, not that their contacts are with
each other.

\begin{remark}[On the U-shaped mortality curve]
\label{sm:rem-ushape}
The earlier chapter derived the U-shaped age-mortality curve of influenza,
high in infants and the elderly, from template matching: the elderly
occupied the state an evolved template fits best. The present framework
offers no such derivation, and I do not supply one. Severity in the sense of
\Cref{sm:def-virulence} is a property of a host network and a medium. Its
dependence on age would have to come from how host networks change with
age, which is a matter for the kinetics of \Cref{c1:in-steady} and not for
contact structure. Pointing at immunosenescence~\citep{goronzy2013understanding} names
the right domain without deriving the curve. The curve is left as an
observation this chapter does not explain.
\end{remark}

%% -------------------------------------------------------------------------
\section{What this replaces}
\label{sm:sec:replaces}
%% -------------------------------------------------------------------------

The chapter this one replaces made four principal claims. Their status is
now:

\begin{enumerate}[nosep,label=(\arabic*)]
\item \emph{Exploitation occurs when a template overlaps an agent state.}
      Replaced. An infection is an added-reaction lesion
      (\Cref{sm:def-infection}), and whether it is pathogenic is decided by
      the viability set, not by an overlap integral. The template formalism
      also sits badly with \Cref{c1:prop-template}, which denies that a
      bounded system can hold a complete description of another.
\item \emph{Coordination regime determines vulnerability, and environmental
      properties do not.} Half kept, half withdrawn. Contact structure moves
      the epidemic threshold through $\langle k^2\rangle/\langle k\rangle$,
      exactly and checkably (\Cref{sm:cor-mean,sm:cor-structure}). The
      environment is not thereby excluded; it enters through $\tau$, and for
      vector-borne disease it does so decisively
      (\Cref{sm:sec:environment}).
\item \emph{Vulnerability rises monotonically with coordination, by a factor
      of five from turbulent to phase-locked.} Refuted. The structural
      quantity is U-shaped in the order parameter, not monotone
      (\Cref{sm:sec:weighting}), and concentration is self-limiting
      (\Cref{sm:prop-nonmonotone}).
\item \emph{Tropical disease prevalence follows from co-evolution time, with
      a tenfold ratio.} Withdrawn as quantitative, kept as direction.
      Selection on added lesions is real (\Cref{sm:prop-coev}); the ratio
      and the regional figures were not measured
      (\Cref{sm:sec:coevolution}).
\end{enumerate}

The general theory of exploitation that framed the earlier chapter---the
claim that ecosystems with invasive species, economies with innovations and
software systems with exploits all obey one overlap integral---is also
dropped. Nothing in this chapter extends beyond hosts with reaction
networks, and I no longer think the generality was earned.

%% -------------------------------------------------------------------------
\section{What this chapter does not say}
\label{sm:sec:limits}
%% -------------------------------------------------------------------------

\paragraph{No disease is modelled end to end.} The chapter connects a
within-host definition to a between-host threshold, and computes neither for
any real pathogen. The within-host half would need the kinetics of
\Cref{c1:in-steady}; the between-host half would need a measured contact
distribution and a measured $\tau$.

\paragraph{The checks are structural, not epidemiological.} Every number in
\texttt{validation/ch13/} is a function of contact structure and a
stipulated $\tau$. The absolute values have no epidemiological content; only
the ratios at fixed mean contact do. Nothing is calibrated against observed
incidence anywhere.

\paragraph{The contact models are idealised.} The configuration model of
check C1 has no clustering and no degree correlations, both present in
measured contact networks and both affecting the threshold. The households
model of check C3 has one level of grouping and fully-mixed groups.
\Cref{sm:eq-r0} itself assumes random mixing subject to the degree
distribution, which is exactly the assumption that \Cref{sm:sec:age}
questions.

\paragraph{Immunity is absent.} The reproduction number of
\Cref{sm:thm-second-moment} is defined in a fully susceptible population.
Recovery, waning and vaccination all change the susceptible set, and none
appears here. In the language of \Cref{ch:what-disease-is}, immunity is a
change to the host network that alters whether an arriving $P$ can carry
flux---an (L1) or (L3) change interacting with an (L2) one---and that
interaction is not worked out.

\paragraph{Evolution is named, not modelled.} \Cref{sm:in-selection} is
supplied and \Cref{sm:prop-coev} only says what selection may act on. No
dynamics, no rates, and no prediction about any particular pathogen's
trajectory follow from this chapter.

%% -------------------------------------------------------------------------
\section{The chain, link by link}
\label{sm:sec:ledger}
%% -------------------------------------------------------------------------

\begin{table}[!t]
\centering
\small
\caption{The chain of this chapter. \lnkD{} derived here; \lnkS{} supplied,
with a reference; \lnkC{} checked by computation in
\texttt{validation/ch13/}. The \lnkC{} labels here certify properties of
contact structure, not of any disease.}
\label{sm:tab-ledger}
\begin{tabular}{@{}p{4.4cm}p{8.6cm}p{2.1cm}@{}}
\toprule
Link & Content & Status\\
\midrule
Infection & a pathogen is an added-reaction lesion (L2) & \lnkD\\
No holonomy change & an infection cannot alter any retained cycle's
  holonomy & \lnkD\\
Virulence & severity of that lesion; a relation of agent, host and medium
  & \lnkD\\
Host range & reactions that cannot run lie on no cycle, so add no lesion
  & \lnkD\\
Threshold & $R_0=\tau\langle k^2\rangle/\langle k\rangle$ & \lnkS\\
Mean is not enough & equal mean contact, unbounded ratio of $R_0$
  & \lnkD\lnkC\\
Structure alone & contact change crosses the threshold at fixed $\tau$
  & \lnkD\lnkC\\
Group size & $R_*$ non-monotone; total concentration transmits nothing
  & \lnkD\lnkC\\
Weighting function & U-shaped in the order parameter, not monotone;
  dropped & \lnkC\\
Environment & enters through $\tau$, through contact, or through the host
  & \lnkD\\
Malaria thermal curve & transmission peaks in the mid-twenties~$^{\circ}$C;
  highlands excluded & \lnkS\\
Selection & acts on which reactions are added and their conductances
  & \lnkD\lnkS\\
Co-evolution ratio & withdrawn; direction kept, no number & ---\\
Age-assortative contact & measured in diary studies & \lnkS\\
Assortativity alone & no threshold effect at constant row sums
  & \lnkD\lnkC\\
Row-sum inequality & raises the threshold quantity at constant mean
  & \lnkD\lnkC\\
Latency, U-shape & phenomena named; mechanisms not derived here
  & \lnkS\\
\bottomrule
\end{tabular}
\end{table}

\Cref{sm:tab-ledger} collects the labels. The pattern differs from
\Cref{c1:tab-ledger} in a way worth noticing. In
\Cref{ch:what-disease-is} the derived links ran all the way to a real cell,
because stoichiometry was enough to reach it. Here the derived links reach a
threshold formula and stop, and everything beyond---which diseases are
vector-borne, how contact is distributed in a real population, what
immunosenescence does---is supplied or absent. That is the honest shape of
the subject: the step from one host to many introduces quantities that no
amount of derivation from \Cref{c1:post-counted} will produce.
